\documentclass[11pt]{article}
\usepackage{amsmath,amssymb,amsthm}
\newtheorem{proposition}{Proposition}
\title{Protected tangent scaling and the exact mixed-entry response threshold}\author{}\date{}
\begin{document}\maketitle
\paragraph{Target and distinction from the elementary stencil estimate.}
Keep the clean-born stored $d=2$ stable $H$ packet of the clean born first entry bounds calculation. No original fine
primitive is restored. The inverse assigned-length factor in the elementary
retained-coordinate stencil estimate is an upper-bound loss, not an intrinsic
scaling law for a physical tangent. The relevant tangent is to the SAME
completed determining problem that supplies its stored physical representative.
This note differentiates the manuscript's exact DA scaling on such a problem,
then states the quantitative mixed-entry assertion still missing.

\begin{proposition}[Differentiated protected scaling]
Let the actual common field $U(v)=\exp(a A(v))$ obey DA1 on the complete protected
cube containing $x_0+rX$ and its collars. Suppose a current-coordinate direction
$w$ has the actual physical response
\[
 \|D_vA[w]\|_{C^{1,\beta}}+
 \|D_vJ_a(U)[w]\|_\infty\le C_{\rm resp}g_e\|w\|,
 \quad \|A\|_{C^{1,\beta}}+\|J_a(U)\|_\infty\le\epsilon.       \tag{1}
\]
The constant retains the original determining graph and carrier geometry and
is independent of $a/r$. Use the DA gauge $g(v,x)=\exp S(v,x)$,
$S=\sum_i(x_i-x_{0,i})A_i(v,x_0)$, and old coordinates
$y=(x-x_0)/r$, $\xi=a/r$. Then, on the same strict chart margin,
\[
 \|D_v\widetilde A[w]\|_{C^{1,\beta}_\xi}
       \le C_D C_{\rm resp} r^2 g_e\|w\|,
 \qquad
 \|D_vJ_\xi(U^g)[w]\|_\infty
       \le C_D C_{\rm resp} r^3 g_e\|w\|.              \tag{2}
\]
\end{proposition}
\begin{proof}
The logarithmic derivative $K=(D_vg[w])g^{-1}$ is an integral of conjugates
of $D_vS[w]$, so $\|K\|\le C_D r C_{\rm resp}g_e\|w\|$ and its first
current-coordinate spatial differences are bounded without an inverse fine mesh.
Differentiating DA4 preserves its exact zero at $x_0$. Its linear part is the
response difference $D_vA_i(x)-D_vA_i(x_0)$, bounded by
$C_D r C_{\rm resp}g_e\|w\|$; the analytic remainder and its first/H\"older
spatial differences obey the same differentiated DA3 bounds on the common
chart. Rescaling $\widetilde A=r A^g$ gives the first estimate in (2).

For the second, use the EXACT Wilson-current identity DA6, not a generic
second-difference estimate:
\[
 D_vJ_\xi(U^g)[w]
 =r^3\{\operatorname{Ad}_g D_vJ_a(U)[w]
                 +[K,\operatorname{Ad}_gJ_a(U)]\}.      \tag{3}
\]
The first term uses the physical-current part of (1); the second uses the
bounds on $K,J_a$ and the controlled adjoint norm. Both have the stated bound.
No estimate of a second fine difference by $a^{-1}$ has been used.
\end{proof}

\paragraph{Stored-packet Ward cancellation.}
For the complete gauge-invariant physical representative of $R_j$, differentiating
$R_j(U^g,J_{\rm phys}(U^g))=R_j(U,J_{\rm phys}(U))$ shows that the simultaneous
pure gauge tangent of all links, current legs and transports has zero derivative.
This is the actual gauge Ward identity for the stored function, not a claim that
an independent-current partial derivative vanishes. Saved covariant completions,
anchor transports and projection partners must transform together. It justifies
using the analysis gauge in (2) without a new gauge derivative cost in that
complete scalar. It neither deletes a whole-kernel selector boundary response nor
identifies two different determining problems.

\paragraph{Required exponent for the actual mixed response.}
For use in the CANONICAL scalar envelope, a tangent estimate alone is not
sufficient. One must also have the actual zero field/current at flat determining
data, and a common admitted scaling path for ALL determining, exterior and
pinned inputs (MT22, RG6/10). Otherwise the nonzero physical baseline remains
explicit and may cost the physical norm in the clean born first entry bounds calculation. Only with that additional
zero/path condition does integration of the tangent bound give the canonical
amplitude used in RG31.

For the literal first-entry stored readout on that canonical family, a sufficient response is (1) with
$C_{\rm resp}$ uniform in depth and fine mesh, up to the already paid carrier
polynomials and admissible relative coupling factors. The field/current radius
conditions remain $C_Dr^2\epsilon< c_{\rm ad}\alpha_j$ and
$C_Dr^3\epsilon<c_{\rm ad}\gamma_j$; there must be a strict common margin on
all preparation/source walls. This is an exponent-zero response requirement,
not an arbitrary bounded coordinate Jacobian.

More generally, if the actual canonical control incurs $r^{-\eta}$, substitution
into the established RG31 envelope changes its three strict-age exponents to
\[
 s-2\eta,\qquad \beta-3\eta,\qquad 2-4\eta.             \tag{4}
\]
Thus this route still yields SOME strict decay only for
\[
 \eta<\min\{s/2,\beta/3,1/2\}.                         \tag{5}
\]
For a specified target $s_*>0$, it suffices to prove
$\eta\le\min\{(s-s_*)/2,(\beta-s_*)/3,(2-s_*)/4\}$.
At exponent zero the three original weights are preserved. The generic
$h^{-1}=r^{-1}$ fallback lies outside (5). These are consequences of the
available upper bounds, not a lower-bound obstruction to a better argument.
The extra weak-coupling factors in the cubic/quartic terms remain as in the clean born first entry bounds calculation;
(5) is the uniform age criterion without spending them as a new age reserve.

\paragraph{Comparison with the manuscript.}
RG13--14 and LM38--41 provide the field/current response in (1) on an ADMITTED
common regular determining problem. The strong Maxwell response, or the actual
KKT multiplier response, is essential: a first-order potential bound alone does
not provide the physical current derivative. Differentiating
$J(P)=Q(P)^*\lambda$ gives
$DJ=(DQ^*)\lambda+Q^*D\lambda$ only for that actual critical constrained field,
with every prefix/boundary row retained. An auxiliary reference or a separately
prepared joining field is not identified with it by this equation.

DA1 requires a complete protected cube and contractible carrier; it is not a
statement across a wild hole. HI3 collapses only the same-action nested hierarchy
with a common boundary class. HI14 explicitly requires a common joint domain;
HI15 presupposes a protected regular output. LM42 retains the already integrated
representative, but LM45--47 does not prove that its mixed prepared range belongs
to the smaller saved marked core. RG13 explicitly excludes an unprotected
rough-collar conclusion. At the physical wall endpoint, R-II (1.50) IS an original constrained Wilson
minimum: the issue is not absence of every critical equation. The missing
identification is between its CHANGED determining/target/boundary record and
the protected current problem required by RG13, together with literal nesting
into the saved old problem. At the independent three-join walls (1.61)--(1.63),
no common-action critical-map identity is supplied. None of the cited results
verifies (1) on that actual simultaneous mixed domain.

Performing source localization before action localization keeps later kernel
walls out of an already admitted source argument. It does not prove the earlier
changed-record endpoint identity or give the missing protected collar. Likewise,
its real coordinate-energy reserve and its canonical/physical norm distinction
remain those of the clean born first entry bounds calculation; the tangent estimate cannot create an unspent
energy record or pay a fixed physical background as a canonical one.

\paragraph{Necessary model lemma for this route.}
For the original selected first-entry record, identify its COMPLETE current-unit
coordinates (including retained old, exterior and pinned rows), its one source
readout and all proper-difference companions. Prove on the original selected
support and simultaneous wall domain that this readout either (i) is a common
admitted critical field to which (1)--(2) apply, with the original collars and
original energy ownership, or (ii) satisfies a direct completed-packet response
bound with loss exponent obeying (5). The second option can use the saved
projections and true Ward cancellations without requiring each intermediate field
to obey the first option, but such a bound is not supplied by the current text.
This is a model estimate to prove, not an additional assertion made here.

\paragraph{Precise stopping point and verification.}
The original R-I selector/range definitions are not used here, so not all hypotheses
of the original preparation theorem are verified.
Even supplying those definitions would not itself prove their smaller marked
range or the response estimate: that is the unresolved mathematical step.
The clean response algebra, old rescaling, gauge derivative cancellation and
threshold (5) are derived above. No actual nonprotected selected
history satisfying the required response has been verified, so no MT36 clause
is promoted to closed. Further moment computations cannot settle this domain
and physical-tangent identification.
\end{document}
